\section*{Cluster A Step 38: Higher-Layer Shadow Uniqueness}

\paragraph{Grade.}
Finite-carrier diagnostic. The computation derives low-energy shadows after scalar breaking and tests a clean stable-composite plus mass requirement. It does not derive content or values.

\paragraph{Carrier.}
The Step-35 residual is rederived as \(12\): \(8\) in \(\texttt{2|3}\) and \(4\) in \(\texttt{4}\).

\paragraph{Shadow Requirement.}
The base shadow requirement asks for a confining unbroken non-abelian subgroup plus completed mass sector. The clean shadow also rejects massive broken gauge generators charged under the confining subgroup.

\paragraph{Result.}
The base shadow leaves all \(12\), confirming that partial breaking of the single-factor family is viable. The clean shadow leaves \(8\), all in \(\texttt{2|3}\). The \(\texttt{4}\) family leaves a residual \(\texttt{3}\) subgroup but also \(6\) confining-charged broken-vector exotics, so it fails the clean shadow.

\paragraph{Verdict.}
\(\texttt{SHADOW\_UNIQUENESS\_WITH\_CONTENT\_LIMIT}\). The structure family is distinguished by the low-energy shadow only conditionally on the introduced clean-separation requirement: no confining-charged broken vectors. The Step-14 reference support is not uniquely selected among the eight content variants.

This uniqueness is CONDITIONAL on the INTRODUCED clean-separation requirement. Step 39 finds this requirement is not supplied by the toy's internal grounding attempt: the "standard confining sector" grounding is circular. Equivalently, Step 41 records clean-separation as \(\Delta_{\mathrm{fact}}=\varnothing\), a recognition-source closure condition, not a law supplied by the construction.

\paragraph{Boundary.}
This is not a fixed-factor-count rule or a size ordering. The old factor-local and scalar-stage currency diagnostics remain rejected.
